\subsection{Indefinite Iteration: The \texttt{while} Loop}

A \texttt{while} loop is Python's direct mechanism for indefinite iteration. The word indefinite does not mean uncontrolled. It means that the number of repetitions is not fixed by the loop syntax itself. The loop continues as long as its condition keeps evaluating as true.

This is different from a \texttt{for} loop over a known traversal source such as a list, tuple, string, or \texttt{range}. A \texttt{while} loop does not ask an iterable for the next value. It repeatedly asks a condition whether the loop should continue.

\begin{quote}
A \texttt{while} loop is a repeated conditional branch with a back-edge to the condition test.
\end{quote}

The programmer is responsible for making sure that something inside or around the loop eventually changes the condition, unless the loop is intentionally meant to run forever.

\subsubsection{Syntax Mechanics and Condition Evaluation Pipelines}

The syntax of a \texttt{while} loop has three visible parts:

\begin{verbatim}
while condition:
    loop_body
\end{verbatim}

The \texttt{while} keyword introduces the repeated decision point. The condition is evaluated before each possible execution of the body. The colon announces the block. The indented statements form the body that executes when the condition is true.

A small example:

\begin{verbatim}
count = 1

while count <= 3:
    print(f"count: {count}")
    count = count + 1

print("loop finished")
\end{verbatim}

Possible output:

\begin{verbatim}
count: 1
count: 2
count: 3
loop finished
\end{verbatim}

The execution route is:

\begin{verbatim}
set count = 1
test count <= 3
    true:
        print count
        increase count
        go back to test count <= 3
    false:
        continue after the loop
\end{verbatim}

The condition is not checked only once. It is checked before every iteration.

The condition expression produces a routing decision. In this example, the comparison \texttt{count <= 3} produces a \texttt{bool} object:

\begin{verbatim}
count = 1
check = count <= 3

print(f"check: {check}")
print(f"type(check): {type(check)}")
print(f"dir(check) sample: {dir(check)[:4]}")
print(f"check.__bool__(): {check.__bool__()}")
\end{verbatim}

Possible output:

\begin{verbatim}
check: True
type(check): <class 'bool'>
dir(check) sample: ['__abs__', '__add__', '__and__', '__bool__']
check.__bool__(): True
\end{verbatim}

The object \texttt{True} tells the loop to enter the body. When the comparison later produces \texttt{False}, the loop body is skipped and execution continues after the loop.

A \texttt{while} condition does not have to be a comparison. Any object can be truth-tested:

\begin{verbatim}
lines = ["Ni!", "Spam!", "Giddy up!"]

while lines:
    line = lines.pop(0)
    print(f"line: {line}")

print("no lines left")
\end{verbatim}

Possible output:

\begin{verbatim}
line: Ni!
line: Spam!
line: Giddy up!
no lines left
\end{verbatim}

Here the condition is the list object \texttt{lines}. A non-empty list is true. An empty list is false.

\begin{verbatim}
lines = ["Ni!", "Spam!"]

print(f"type(lines): {type(lines)}")
print(f"dir(lines) sample: {dir(lines)[:10]}")
print(f"len(lines): {len(lines)}")
print(f"bool(lines): {bool(lines)}")
print(f"has __len__: {'__len__' in dir(lines)}")
print(f"has pop: {'pop' in dir(lines)}")
\end{verbatim}

Possible output:

\begin{verbatim}
type(lines): <class 'list'>
dir(lines) sample: ['__add__', '__class__', '__class_getitem__', '__contains__', '__delattr__', '__delitem__', '__dir__', '__doc__', '__eq__', '__format__']
len(lines): 2
bool(lines): True
has __len__: True
has pop: True
\end{verbatim}

The \texttt{\_\_len\_\_()} method supports length-based truth-value testing. The \texttt{pop()} method removes and returns an element. In this loop, every call to \texttt{pop(0)} makes the list shorter. Eventually the list becomes empty, and the condition becomes false.

The condition pipeline is therefore:

\begin{verbatim}
evaluate condition expression
    -> truth-test resulting object
        -> true: execute loop body
        -> false: skip loop body and exit loop
\end{verbatim}

A \texttt{while} loop may also begin false and never execute its body:

\begin{verbatim}
count = 5

while count < 3:
    print(f"count: {count}")
    count = count + 1

print("body was not entered")
\end{verbatim}

Possible output:

\begin{verbatim}
body was not entered
\end{verbatim}

This is because \texttt{while} is a pre-test loop. The condition is checked before the first iteration.

\subsubsection{Guarding Against Resource Exhaustion: Engineering Manual Exit Conditions and Loop Invariants}

A \texttt{while} loop can run forever if its condition never becomes false. This is not a special Python problem. It is a direct consequence of the control-flow graph. A loop has a back-edge. If nothing breaks the back-edge, execution keeps returning to the same test.

Example of a dangerous loop:

\begin{verbatim}
count = 1

while count <= 3:
    print(f"count: {count}")
\end{verbatim}

The variable \texttt{count} never changes. The condition \texttt{count <= 3} remains true forever. The program keeps printing until it is externally interrupted or until some resource limit is reached.

A safe \texttt{while} loop usually needs a loop invariant and a progress operation.

A loop invariant is a fact that should remain controlled and meaningful while the loop runs. A progress operation is the update that moves the loop toward termination.

\begin{verbatim}
count = 1
limit = 3

while count <= limit:
    print(f"count: {count}, limit: {limit}")
    count = count + 1

print("safe exit")
\end{verbatim}

Possible output:

\begin{verbatim}
count: 1, limit: 3
count: 2, limit: 3
count: 3, limit: 3
safe exit
\end{verbatim}

In this example:

\begin{itemize}
    \item \texttt{limit} stays fixed;
    \item \texttt{count} starts below or equal to the limit;
    \item each iteration increases \texttt{count};
    \item eventually \texttt{count <= limit} becomes false.
\end{itemize}

The invariant is not a magic formal object. It is the programmer's controlled statement about what should remain true while the loop operates.

For this loop, a useful plain-language invariant is:

\begin{quote}
At the start of each iteration, \texttt{count} is the next number to process, and all smaller numbers from the starting point have already been processed.
\end{quote}

The progress operation is:

\begin{verbatim}
count = count + 1
\end{verbatim}

Without that line, the loop has no internal route toward termination.

A second common safety pattern is a manual attempt counter:

\begin{verbatim}
attempt = 1
max_attempts = 3
password_ok = False

while not password_ok and attempt <= max_attempts:
    print(f"attempt {attempt}: checking password")
    attempt = attempt + 1

print("login loop stopped")
\end{verbatim}

Possible output:

\begin{verbatim}
attempt 1: checking password
attempt 2: checking password
attempt 3: checking password
login loop stopped
\end{verbatim}

The condition has two parts:

\begin{verbatim}
not password_ok
attempt <= max_attempts
\end{verbatim}

The loop continues only while the password is not accepted and the attempt limit has not been exceeded.

The state objects are ordinary runtime objects:

\begin{verbatim}
attempt = 1
max_attempts = 3
password_ok = False

print(f"type(attempt): {type(attempt)}")
print(f"dir(attempt) sample: {dir(attempt)[:8]}")

print(f"type(password_ok): {type(password_ok)}")
print(f"dir(password_ok) sample: {dir(password_ok)[:8]}")
print(f"password_ok.__bool__(): {password_ok.__bool__()}")
\end{verbatim}

Possible output:

\begin{verbatim}
type(attempt): <class 'int'>
dir(attempt) sample: ['__abs__', '__add__', '__and__', '__bool__', '__ceil__', '__class__', '__delattr__', '__dir__']
type(password_ok): <class 'bool'>
dir(password_ok) sample: ['__abs__', '__add__', '__and__', '__bool__', '__ceil__', '__class__', '__delattr__', '__dir__']
password_ok.__bool__(): False
\end{verbatim}

The integer object supports arithmetic methods such as \texttt{\_\_add\_\_()}, which is used by \texttt{attempt + 1}. The boolean object supports \texttt{\_\_bool\_\_()}, which participates in the condition.

A third useful safety technique is to place a clear stop condition near the top of the loop body:

\begin{verbatim}
energy = 5

while True:
    print(f"energy: {energy}")

    if energy <= 0:
        print("No soup for the loop.")
        break

    energy = energy - 2

print("loop ended")
\end{verbatim}

Possible output:

\begin{verbatim}
energy: 5
energy: 3
energy: 1
energy: -1
No soup for the loop.
loop ended
\end{verbatim}

This form uses \texttt{while True:} to create an intentionally open loop, then uses \texttt{break} to leave it. The detailed mechanics of \texttt{break} will be studied later. Here, the important point is that even an apparently infinite loop must have a deliberate exit route if it is not meant to run forever.

A safe \texttt{while} loop should usually answer these questions:

\begin{itemize}
    \item What state does the condition inspect?
    \item What state changes during each iteration?
    \item Why must the condition eventually become false?
    \item What prevents the loop from consuming unbounded time, memory, input, or output?
\end{itemize}

For example, this loop has bounded output:

\begin{verbatim}
n = 1

while n <= 4:
    print(f"{n:>2} | {n * 42:>3}")
    n = n + 1

print("table complete")
\end{verbatim}

Possible output:

\begin{verbatim}
 1 |  42
 2 |  84
 3 | 126
 4 | 168
table complete
\end{verbatim}

The formatting is not part of the loop logic, but it helps expose the loop state. The condition controls how many rows can be produced. The update moves the loop to the next row. The fixed limit prevents unbounded printing.

The practical rule is:

\begin{quote}
A \texttt{while} loop is safe when its continuation condition is tied to state that moves predictably toward an exit.
\end{quote}

This is the programmer's responsibility. Python supplies the syntax, the truth-value test, and the back-edge. It does not automatically prove that the back-edge will eventually stop.